\chapter*{\textbf{Introduction générale}}
\addstarredchapter{Introduction générale}
\label{intro}
\indent Dans l'industrie pharmaceutique et dermo-cosmétique, la qualité n'est pas un objectif 
parmi d'autres : elle constitue une exigence fondamentale. Les produits développés étant 
destinés à des patients et des consommateurs, leur qualité et leur sécurité doivent être 
maîtrisées tout au long de leur cycle de vie. Cette exigence s'étend aux partenaires et sous
traitants avec lesquels l'entreprise collabore (désignés sous le terme de Tierces Parties dans 
le référentiel qualité du groupe), dont l'évaluation et le suivi participent directement à la maîtrise 
des risques. Cette culture de la qualité est au cœur des activités du groupe Pierre Fabre, 
laboratoire pharmaceutique et dermo-cosmétique français fondé en 1951 à Castres et 
aujourd'hui présent dans 120 pays. 

\indent Assurer le suivi de ces Tierces Parties à grande échelle nécessite de disposer d'une 
donnée fiable, structurée et exploitable. C'est autour de cet enjeu que s'articule IPANEMA. 
Les informations relatives aux Tierces Parties sont collectées à travers une application Power 
Apps et stockées dans différentes listes SharePoint. Jusqu'alors, elles étaient consolidées au 
moyen d'un workflow Alteryx en une table unique, puis restituées dans un dashboard Tableau. 
Cette architecture a cependant révélé plusieurs limites : lignes dupliquées ou manquantes lors 
des jointures, incohérences dans certains calculs et maintenance complexe dans un 
environnement progressivement orienté vers l'écosystème Microsoft.

\indent Le projet a donc consisté à repenser la chaîne de traitement et de valorisation des 
données d'IPANEMA, en construisant un modèle de données adapté et en migrant la partie 
décisionnelle vers Power BI. L'objectif était de proposer aux équipes Qualité une solution plus 
fiable, maintenable et intégrée au système d'information du groupe, leur permettant de 
disposer d'une vision cohérente des informations nécessaires au suivi des Tierces Parties. Au
delà de l'aspect technique, la fiabilisation de ces données contribue ainsi à une meilleure 
maîtrise de la qualité tout au long de la chaîne de valeur.

\indent Ce projet s'inscrit dans ma première année d'alternance au sein de la R&D Dermo
Cosmétique & Personal Care (DCPC) de Pierre Fabre, à Toulouse-Langlade, dans le cadre 
d'un projet démarré en février 2026. En tant que Data Analyst au sein de l'équipe Digital 
Product Dev et Perf, j'ai assuré la conduite du projet IPANEMA dans son ensemble, sous le 
cadrage de mon maître d'apprentissage, depuis le recueil et l'analyse des besoins jusqu'au 
développement et au déploiement de la solution auprès des utilisateurs finaux. Cette double 
dimension technique et pilotage m'a placée à l'interface entre les besoins métier des équipes 
Qualité et les enjeux techniques liés à la structuration, au traitement et à la valorisation des 
données.

\indent Dès lors, la problématique qui guidera ce mémoire est la suivante : comment repenser 
la structuration et la restitution des données d'IPANEMA afin de fiabiliser le suivi des Tierces 
Parties, tout en proposant aux équipes Qualité un outil de pilotage performant, maintenable et 
intégré à l'écosystème Microsoft du groupe ?

\indent Pour répondre à cette problématique, ce mémoire s'articule autour de quatre parties. 
La première, le rapport d'action, présente une synthèse du projet, de ses enjeux et des 
résultats obtenus. La deuxième, le rapport d'étude, expose l'analyse de l'existant, la démarche 
adoptée ainsi que les choix fonctionnels et techniques ayant conduit à la solution retenue. La 
troisième partie revient sur le management du projet, son organisation, sa planification et son 
pilotage. Enfin, la dernière partie propose un bilan personnel et professionnel, mettant en 
perspective les compétences développées et les enseignements tirés de cette première année 
d'alternance.